%! The Language of Variation — Unit 1, Lesson 1.1 — AP Practice
% GRADUAL RELEASE: AP-format practice, assigned but NOT scored — the cover's score column
% reads NA for this component. The graded practice is the `homework` component that follows it
% at the back of the packet. Shaped like the AP exam: Section I multiple choice in
% context, Section II free response with an interpret-in-context part.
% Part (d) of the free response is the spiral item — it reaches back to Lesson 1.0's
% statistical-question test and makes students run it on a categorical variable.
% Page lockstep with ap_practice_key rests on \writelines{n} in the blank matching a terse
% \ansline{} in the key — keep key answers short enough not to wrap past the reserved lines.
% The next-lesson preview is NOT here: it closes the packet, in ../homework/.
\documentclass[12pt]{article}
\usepackage{apstats-article}
\usepackage{apstats-boxes}

\begin{document}

\pageheader{Unit 1, Lesson 1.1}{AP Practice}
% NAMESTRIP: no \namedateperiod here — the name row lives on the cover only.

\vspace{0.15in}

% Authored BYTE-IDENTICALLY in ap_practice/ and ap_practice_key/ — this box is what keeps the
% blank and the key the same number of pages.
\boxguard[8]
\begin{remindbox}
\small
This page is \textbf{not required and not scored} --- your graded homework is the last section
of this packet. These are AP-format reps: work them for the practice, and bring what you get
stuck on to office hours or study hall.
\end{remindbox}

\vspace{0.12in}

\boxguard[30]
\begin{scenariobox}[The Trailhead Register]{navy}
\small
Every party that hikes at Ross Hollow State Park signs the register at the trailhead. Here is
last season's, with the columns the ranger added:

\vspace{8pt}
\renewcommand{\arraystretch}{1.3}
\begin{tabularx}{\linewidth}{@{} l Y Y Y Y Y Y @{}}
\rowcolor{sky}
\textbf{Party} & \textbf{Trail} & \textbf{Group size} & \textbf{Time out (min)} &
\textbf{Dog along} & \textbf{Permit \#} & \textbf{Home ZIP} \\
\hline
Alvarez  & Ridge Loop   & 4 &  95 & Yes & 2041 & 80211 \\
Boone    & Cedar Falls  & 2 & 140 & No  & 6318 & 80304 \\
Chen     & Ridge Loop   & 6 &  80 & Yes & 4772 & 81657 \\
Delgado  & Summit Trail & 3 & 215 & No  & 1150 & 80211 \\
Egan     & Cedar Falls  & 2 & 155 & Yes & 3907 & 80304 \\
Fontaine & Summit Trail & 5 & 190 & No  & 5284 & 80211 \\
\hline
\end{tabularx}
\end{scenariobox}

\vspace{0.14in}

\boxguard[16]
\begin{headlinebox}{goldbg}
\small
\textbf{Section I --- Multiple Choice.} Circle the single best answer. Every answer choice is
about the context, not just the arithmetic --- read them all before choosing.
\end{headlinebox}

\vspace{0.10in}

\begin{enumerate}[leftmargin=1.9em, itemsep=0.14in]

  \item Which column of the register above records a \textbf{quantitative} variable?
        \begin{enumerate}[label=\textbf{(\Alph*)}, leftmargin=2.2em, itemsep=3pt]
          \item Trail
          \item Dog along
          \item Time out (min)
          \item Permit \#
          \item Home ZIP
        \end{enumerate}

  \item A new ranger argues that \emph{Home ZIP} must be quantitative ``because every entry is
        a five-digit number.'' The best response is that
        \begin{enumerate}[label=\textbf{(\Alph*)}, leftmargin=2.2em, itemsep=3pt]
          \item ZIP codes are quantitative, so the ranger is correct.
          \item ZIP codes are categorical, because their digits identify an area rather than
                record an amount.
          \item ZIP codes are categorical, because five digits is too many to average.
          \item ZIP codes are quantitative, because they can be sorted from smallest to largest.
          \item the classification depends on how many parties signed the register.
        \end{enumerate}

  \item The ranger's report states, ``Half the parties in this register hiked with a dog.''
        This statement is
        \begin{enumerate}[label=\textbf{(\Alph*)}, leftmargin=2.2em, itemsep=3pt]
          \item evidence that \emph{Dog along} is a quantitative variable, because it produced
                a number.
          \item evidence that \emph{Dog along} is quantitative, because percents are numerical.
          \item a numerical summary of a categorical variable, which remains categorical.
          \item impossible, because a categorical variable cannot be summarized.
          \item a numerical summary of a quantitative variable.
        \end{enumerate}

  \item Suppose the ranger stops recording \emph{Time out} in minutes and instead records each
        hike as \emph{Short}, \emph{Medium}, or \emph{Long}. Which statement is \textbf{true}?
        \begin{enumerate}[label=\textbf{(\Alph*)}, leftmargin=2.2em, itemsep=3pt]
          \item The variable stays quantitative, because the categories came from minutes.
          \item The variable becomes categorical, and the mean time on the trail can no longer
                be computed from it.
          \item The variable becomes categorical, but the mean time can still be computed from
                the labels.
          \item Neither version of the variable is quantitative.
          \item The change affects the individuals rather than the variable.
        \end{enumerate}

\end{enumerate}

\vspace{0.14in}

\boxguard[16]
\begin{headlinebox}{goldbg}
\small
\textbf{Section II --- Free Response.} Show your reasoning. Every answer must be written
\emph{in the context of the problem} --- that is what earns credit on the AP exam.
\end{headlinebox}

\vspace{0.10in}

\begin{enumerate}[leftmargin=1.9em, itemsep=0.16in, resume]

  \item A city library records four things about each of its 900 cardholders: the
        \textbf{card number}, the \textbf{number of items} checked out last year, the
        \textbf{branch} used most often, and the cardholder's \textbf{age} in years.

        \vspace{6pt}
        \textbf{(a)} Identify the individuals, then classify each of the four variables as
        categorical or quantitative.\\[4pt]
        \writelines{3}

        \vspace{6pt}
        \textbf{(b)} A staff member reports that the average card number is 44{,}318 and calls
        it ``a useful summary of our cardholders.'' Explain why it is not.\\[4pt]
        \writelines{2}

        \vspace{6pt}
        \textbf{(c)} The library replaces \emph{age in years} with \emph{age band}
        (child / teen / adult / senior). Explain how this changes the kind of variable, and
        name one thing the library can no longer compute.\\[4pt]
        \writelines{3}

        \vspace{6pt}
        \textbf{(d)} Write one statistical question the library could answer using a
        \textbf{categorical} variable from this data set, and state what varies in it.\\[4pt]
        \writelines{2}

\end{enumerate}


\end{document}
